\documentclass[11pt, letterpaper]{article}
\usepackage[utf8]{inputenc}
\usepackage[spanish]{babel}
\usepackage{amsmath, amssymb, amsthm}
\usepackage{geometry}
\geometry{a4paper, margin=1in}

% Definición de entornos para Teoremas y Demostraciones
\newtheorem{theorem}{Teorema}[section]
\newtheorem{exercise}{Ejercicio}
\renewcommand{\qedsymbol}{$\blacksquare$} % Cuadro negro al terminar la demostración

\title{\textbf{Resolución de Ejercicios}\\ \large How to Prove It / Complex Variables}
\author{LuisAngel - Centro de Investigación en Ciencias (CInC)}
\date{\today}

\begin{document}

\maketitle

\section*{Capítulo 1: Lógica Proposicional y Conjuntos}

\begin{exercise}
Sea $A$ y $B$ conjuntos. Demostrar que si $A \subseteq B$, entonces $A \cup B = B$.
\end{exercise}

\begin{proof}
Supongamos que $A \subseteq B$. Debemos demostrar que $A \cup B = B$. PUTODFAFSAFASFSAFDSAFSAFASFSADFDAS
Por definición de igualdad de conjuntos, esto requiere probar dos cosas:DFAFSAFASFSAFDSAFSAFASFSADFDAS
\begin{enumerate}
    \item $A \cup B \subseteq B$
    \item $B \subseteq A \cup B$
\end{enumerate}
(Aquí va el resto de tu desarrollo matemático usando $x \in A \implies x \in B$).
\end{proof}

\section*{Capítulo 2: Variables Complejas (Churchill)}
% Usa \mathbb{C} para los números complejos
Para cualquier $z \in \mathbb{C}$ donde $z = x + iy$, el módulo se define como $|z| = \sqrt{x^2 + y^2}$.

\end{document}
